We next assessed the sensitivity of the key associations while retaining
the final mixed beta specification as the reference model.
Table~\ref{tab:robustness_final} reports standardized coefficients for Bias
Discrepancy and Toxicity after omitting Content Reaction Positivity,
restricting pages by the number of distinct bias-scored links in their
observed sharing histories, restricting posts to the Twitter collection
window (October 1--December 11, 2018), reconstructing page bias without the
focal link, and varying the minimum-reaction criterion.

For the sharing-history restrictions, we counted only distinct links with
available ideological scores, which are the links that contribute to $b_p$.
The page-activity control retained its original definition based on all
distinct shared links. In the leave-link-out analysis, we removed all posts
by a page sharing the focal news link from that page's full observed history,
then recalculated page leaning, page extremity, and Bias Discrepancy together.
This excluded 371 posts from pages with no remaining bias-scored history,
leaving 4,213 posts; an all-centrist remaining history was assigned $b_p=0$,
consistent with the original calculation. Other predictors retained their
original definitions. The 25-reaction sample was rebuilt from the source
data using the same preprocessing and metric-construction procedures before
applying the lower engagement cutoff. In every fitted sample, continuous
predictors were standardized and the boundary adjustment of $RHI$ was
recomputed using that sample's size.

\begin{table}[htb]
\centering
\footnotesize
\setlength{\tabcolsep}{3pt}
\caption{Sensitivity analyses using the final mixed beta model for the
Reaction Homogeneity Index. Except for the stated modification, each model
retains the fixed- and random-effects specification in
Table~\ref{tab:beta_results}. Continuous predictors are standardized within
each fitted sample. Coefficients are on the logit scale and intervals are
two-sided 95\% Wald confidence intervals.}
\label{tab:robustness_final}
\begin{tabular}{p{3.7cm}ccccc}
\toprule
\textbf{Specification} & $\beta_{\Delta_b}$ & \textbf{95\% CI}
& $\beta_{\mathrm{Tox}}$ & \textbf{95\% CI} & \textbf{N} \\
\midrule
Final model & $-0.173$ & $[-0.211, -0.135]$ & $-0.149$ & $[-0.196, -0.102]$ & 4,584 \\
CRP omitted & $-0.178$ & $[-0.219, -0.138]$ & $-0.133$ & $[-0.180, -0.086]$ & 4,584 \\
Pages with $\geq 2$ bias-scored links & $-0.201$ & $[-0.241, -0.160]$ & $-0.162$ & $[-0.211, -0.113]$ & 4,213 \\
Pages with $\geq 5$ bias-scored links & $-0.219$ & $[-0.264, -0.175]$ & $-0.170$ & $[-0.226, -0.113]$ & 3,405 \\
Twitter-window posts only & $-0.173$ & $[-0.214, -0.132]$ & $-0.144$ & $[-0.192, -0.096]$ & 4,099 \\
Leave-link-out page bias & $-0.192$ & $[-0.234, -0.150]$ & $-0.165$ & $[-0.214, -0.116]$ & 4,213 \\
Minimum 25 reactions & $-0.171$ & $[-0.206, -0.135]$ & $-0.133$ & $[-0.179, -0.088]$ & 5,774 \\
Minimum 100 reactions & $-0.167$ & $[-0.209, -0.126]$ & $-0.177$ & $[-0.229, -0.126]$ & 3,558 \\
\bottomrule
\end{tabular}
\end{table}

Across all eight specifications, Bias Discrepancy and Toxicity had
negative coefficients, with all 95\% confidence intervals excluding zero.
The Bias Discrepancy association persisted after removing the focal link
from the page-bias calculation, after restricting the sample to posts
within the Twitter measurement window, and under alternative engagement
cutoffs. These checks support the robustness of the reported association
to these specific analytical choices. They do not remove the limitations
of estimating page ideology from observed sharing behavior or establish a
causal effect. Because standardization was performed within each sample,
coefficient magnitudes reflect the predictor distributions in those samples.

Omitting $CRP$ from the final specification reduced marginal $R^2$
from $0.762$ to $0.154$, while conditional
$R^2$ remained similar ($0.977$ versus $0.979$).
Marginal and conditional $R^2$ were calculated using the Nakagawa
method implemented in \texttt{performance::r2\_nakagawa()}, with the
lognormal approximation. This indicates that $CRP$ accounts for substantial
fixed-effect variation, while the Bias Discrepancy association persists
when this same-link reaction covariate is omitted.
